\section{Source-bound certificates for binary normal-surface topology}
\label{sec:normal-certificates}

The next report-47 integration connects the local orbit proofs of
Section~\ref{sec:orbit-certificates} to the maintained normal-surface adapter.
The report's polygon-stack adapter is useful as an independent reference,
but leaves ambient-manifold validity as a precondition. The maintained
adapter already checks the finite manifold itself. We retain its edge-point
interval model and manifold validator, and add optional certificates to its
three existing orbit queries. Section~\ref{sec:normal-multiplicity} subsequently
reduces common coordinate multiplicity and introduces a compatible second
proof version; the measurements here retain the direct version-one path.
Thus this integration strengthens the
checkable evidence without weakening the existing ambient input contract.

\subsection{Exact scope and reconstructed geometry}

The input is a reciprocal tetrahedral face-pairing array and seven
nonnegative normal coordinates per tetrahedron. The native validator checks
connectedness, ambient orientability, absence of reversed-edge
identifications, circle or interval edge links, sphere or disc vertex links,
and exactly one finite torus boundary. Matching equations, quadrilateral
constraints and globally consistent edge-intersection counts are checked
from the supplied coordinates. No recognition of a knot exterior, and no
relation to an input planar diagram, is inferred from these checks.

For an ambient edge $e$, the normal intersection weight is $w_e$.
The surface's one-skeleton is represented on the disjoint union of these
edge-point intervals. At a face corner, the endpoints of its normal arcs
are paired by their rank from the corner. A parity union--find determines
which end of an edge interval is nearest the corner; vertex labels alone
would be insufficient in a one-vertex triangulation. The already validated
normal polygons attach to this one-skeleton, so its connected components
are exactly those of the surface. Restricting to boundary faces gives its
boundary circles. Neither construction enumerates individual discs,
intersection points or sheets.

The finite geometry routines now live in
\path{../fast/fastunknot/normal_surface_geometry.py}; their validation and
interval construction are shared by producer and checker. This is an
explicit common trust boundary. The checker does not invoke the topology
producer, orbit scheduler or boundary-cohomology search. Independent
polygon-stack extraction and Regina supply cross-checks of the shared
geometry rather than being required production dependencies.

\subsection{The three orbit proofs and orientation doubling}

The optional \texttt{normal-surface-topology-v1} certificate contains the
normalized source digest, the complete proofs for \texttt{surface},
\texttt{double} and \texttt{boundary}, a boundary-cohomology witness and a
precisely typed table of topological claims. Integer coordinates and their
hexadecimal encodings have the same normalized source digest. The exact
face records remain bound; an isomorphic relabelling requires fresh
reconstruction and evidence. The digest is a source consistency check, not
a replacement for mathematical verification.

Replay revalidates the caller's manifold and coordinates and reconstructs
all three ordered interval systems. It then invokes the independent local
orbit verifier on those reconstructed inputs. Rewriting the outer digest
cannot make a foreign trace acceptable. Each orbit trace must finish;
an otherwise valid prefix supplies no count.

Let $C$, $C_2$ and $b$ be the verified counts for the normal vector $x$,
its double $2x$, and its boundary. The doubled normal surface is the
horizontal boundary of a regular normal interval neighbourhood of the
original. In an orientable ambient manifold this is its orientation double
cover. An orientable component contributes two lifted components and a
nonorientable component contributes one. Hence
\[
 C_{\mathrm{or}}=C_2-C,\qquad C_{\mathrm{nonor}}=2C-C_2.
\]
The checker rejects negative values. The finite reconstruction also computes
$\chi=V-E+F$: $V=\sum_e w_e$, $E$ counts normal arcs after face
identifications, and $F$ is the sum of normal coordinates. All these are
binary arithmetic operations on a compressed description.

For a connected orientable surface, the checker verifies
$g=(2-b-\chi)/2$, including nonnegativity and integrality. For a connected
nonorientable surface it verifies the positive crosscap number
$2-b-\chi$. Every supplied component count, Euler value, disc count,
genus or crosscap claim must match the reconstructed calculation. Boolean
values are not accepted as integer counts. The empty surface has zero
components and no connected-surface genus claim.

\subsection{A finite witness for essential boundary}

Normal intersections with each ambient boundary edge define a mod-two
cochain $z(e)=w_e\bmod2$. The normal matching equations make it a cocycle:
on each boundary triangle the endpoints of normal arcs occur in pairs.
Its cohomology class is dual to the surface boundary's homology class.
The graph used below has one edge for each ambient boundary edge, retaining
loops and multiple edges, and is independent of sheet multiplicity.

The report supplies two complementary finite witnesses. To prove that the
class vanishes, the producer supplies one bit $p(v)$ per graph vertex, and
the checker verifies
\[
 z(uv)=p(u)+p(v)\pmod2
\]
for every edge. To prove nonvanishing, it supplies a set of graph edges
having even incidence at every vertex and odd sum of $z(e)$. This is a
closed mod-two chain with pairing one against the cocycle. Either witness
has size linear in the finite boundary graph. The positive witness can be
found by a breadth-first potential assignment: an inconsistent edge closes
an odd cycle using the two parent paths. Replay checks the chain directly
and does not repeat that search.

A connected normal surface with $\chi=1$ and one boundary component is a
disc. Its boundary is a simple closed curve in the validated boundary
torus. Such a curve is essential exactly when its mod-two homology class
is nonzero: every essential simple torus curve represents a primitive
integral class. The checker therefore certifies a compressing disc exactly
when $C=1$, $\chi=1$, $b=1$, and the nonzero boundary witness verifies.
A vertex-linking disc correctly fails the last condition. An aggregate
zero boundary class on a disconnected surface says nothing about each
individual boundary curve; the disc conclusion requires connectedness.

\subsection{Complexity, budgets and API}

For $t$ tetrahedra and largest coordinate bit length $B$, each interval
system has $O(t)$ pairings and $O(B+\log(t+1))$-bit endpoints. The ordinary
normal-disc count may be exponential in $B$, but neither the finite
validation nor any of the three orbit queries expands it. The polynomial
AHT production and trace-size bounds of
Section~\ref{sec:native-interval-orbits} therefore yield polynomial-size
normal-topology certificates and polynomial verification in the supplied
encoded geometry and proof length. This is a theorem for a supplied normal
vector, not a polynomial algorithm for finding a compressing disc.

\texttt{normal\_surface\_topology(..., record\_certificate=True)} records the
proofs; the default remains count-only. A single \texttt{max\_cycles}
allowance is shared across all three production queries. An incomplete
query returns \texttt{INCONCLUSIVE} without any topology certificate or
compressing-disc answer. Recording preserves cycle counts and all previous
result fields. Proof production and independent verification are separate
API calls; requesting a proof does not implicitly replay it.

\texttt{verify\_normal\_surface\_certificate} optionally checks a shared
\texttt{max\_operations} cap on the sum of recorded orbit events before
replay. Exceeding it means unverified, not a negative topological claim.
Caller cancellation exceptions propagate, including \texttt{ValueError}.
Invalid supplied manifold data raises the same \texttt{NormalOrbitError}
as production; malformed proof data returns false. Callers can additionally
bound serialized size or provide cooperative deadlines. These controls do
not assert that the entire untrusted input was bounded before allocation.

\subsection{Native validation and performance evidence}

Four delivered source and test files were checked against report 47's
manifest. The native adapter retains the prior finite geometry, whose
count-only results are compared against the actual maintained implementation
at commit \texttt{98f21ce333bf}. Twenty-eight focused tests pass in
11.167 seconds. New proof tests disable all producer search functions during
replay, alter every topological claim, substitute foreign queries, rewrite
source digests, corrupt both boundary-witness forms, exhaust shared caps,
and inject cancellation at early, middle and late checkpoints. Empty,
one-sided, vertex-linking and parallel surfaces include coordinates with
more than 20,000 bits and hexadecimal JSON round trips.

The seeded native audit starts with 252 actual normal vectors in layered solid
tori with one through eight tetrahedra, including compatible sums and empty
vectors. Every result agrees with both Regina and the report's separate
polygon-stack extraction; both native and report certificates replay.
Eighteen further cases attach tetrahedral balls along boundary faces, preserving
the manifold while adding boundary vertices; meridian extensions and each
vertex-linking disc exercise the larger boundary graph. There are twelve
compressing-disc cases, 34 cases with a nonorientable
component and 94 disconnected cases among the 270 total comparisons. The old native count-only result also
agrees exactly in every case. Raw vectors, proof digests and source hashes
are retained in \path{data/normal-certificate-native-audit.json}.

The full maintained suite passes 924 tests with Regina in 246.517 seconds.
The final focused run includes additional odd-cycle and exact-potential
checks on a three-vertex graph. The ten moved geometry definitions are
AST-identical to their baseline versions; the refactoring adds no alternate
manifold-validation path.

The benchmark uses the actual baseline module from that commit with unchanged
orbit-kernel dependencies. Its six arms are the old count-only call, the new
count-only call, an identical new control, optional proof recording,
production followed by independent replay, and replay of an already supplied
proof. There is one shuffled warmup and five shuffled measured rounds per
case. Small fixtures use batches to reduce timer overhead; each call starts
with fresh source validation. Proof transport and JSON serialization are
outside the timed region. All results and produced certificates are compared
outside timing; the combined production/replay arm checks successful replay
inside its call. Source hashes are checked before and after the experiment.

\begin{center}
\begin{tabular}{@{}lrrrrr@{}}
\toprule
Input & count ms & record ms & both ms & replay ms & ratio \\
\midrule
meridian1 & 0.928 & 1.120 & 1.890 & 0.750 & 1.270 \\
meridian8 & 5.043 & 5.662 & 8.596 & 2.453 & 2.079 \\
meridian32 & 70.736 & 82.329 & 95.175 & 16.466 & 4.364 \\
meridian96 & 1267.563 & 1688.191 & 1514.220 & 125.874 & 9.852 \\
parallel5000 & 16.594 & 19.522 & 27.104 & 7.130 & 2.449 \\
one\_sided20000 & 0.740 & 1.164 & 1.624 & 0.728 & 1.059 \\
\bottomrule
\end{tabular}
\end{center}


Here ``both'' means production plus independent verification, and the ratio
is the median paired count-only/replay time. Replay therefore has a different
input contract: a complete proof is already available. The table does not
claim that the full task of finding and verifying a normal surface has been
accelerated by the same ratio. Raw shuffled samples, exact representative
certificates and their encoded sizes are in
\path{../fast/results/normal_certificates_20261008.json}.

The 840 measured calls and 168 excluded warmup calls finish in a
56.791-second run. For the 96-tetrahedron meridian, fresh topology computation
has median 1,267.563 ms and source-bound replay 125.874 ms, with median
paired ratio 9.852. Its proof has 24,038 recorded orbit events and occupies
2,018,995 bytes in compact JSON. This establishes substantially cheaper
checking when a proof is already supplied. Production plus replay takes
1,514.220 ms and is more expensive than count-only discovery.

At 32 tetrahedra, count-only/replay medians are 70.736/16.466 ms (paired
ratio 4.364); at eight tetrahedra they are 5.043/2.453 ms (ratio 2.079).
The $2^{5000}$ parallel multiplicity has ratio 2.450, whereas the 20,001-bit
one-sided example has ratio 1.059. In the latter case finite reconstruction
and binary arithmetic leave little orbit search to avoid. Large coordinate
bit length by itself does not imply a large replay advantage.

Identical current controls have paired ratios from 0.983 to 1.011.
Nevertheless the large fixture's recording-only median, 1,688.191 ms,
exceeds its combined production/replay median. The raw shuffled samples
retain this timing fluctuation; subtracting those separate medians would
not estimate the cost of verification. Nor do the old/new count-only
measurements establish an optimization of discovery: their geometry and
orbit algorithms are unchanged. The useful performance result here is
proof reuse, with independently checked source reconstruction, and its
benefit depends on how much orbit search the source would otherwise need.

\paragraph{What remains before a diagram recognizer can use this.}
A checked compressing disc implies unknottedness only after the supplied
manifold has been established as the original knot exterior. This adapter
validates a finite manifold with torus boundary but supplies no such
correspondence. It also does not search for a normal vector, bound the
cost of a hierarchy that finds one, identify slopes of general components,
or distinguish every component in an aggregate summary. Report 47's
boundary-incidence extensions and the maintained marked-union proofs can
support further work, but neither replaces those missing source and
complexity arguments. General quasi-polynomial recognition is still open
for the maintained implementation.
